% SPDX-FileCopyrightText: 2025 IObundle
%
% SPDX-License-Identifier: MIT

Versat only generates an accelerator's Verilog and software, as shown in Section~\ref{sec:simple_example_outputs}. Turning that output into a SoC peripheral is the job of the \texttt{CreateVersatClass} function in \texttt{iob\_versat.py}, at the root of this repository. The function currently targets the Python setup flow of the IOb-SoC framework, whose cores are classes built on the class-method interface of IOb-SoC's \texttt{iob\_module} (\texttt{\_create\_submodules\_list}, \texttt{\_setup\_confs}, \texttt{\_setup\_ios}, \texttt{\_setup\_regs}). Py2HWSW's current \texttt{iob\_module} is a different, object-based interface without these methods, so a Py2HWSW SoC cannot use \texttt{CreateVersatClass} as it stands. Integrating Versat into Py2HWSW requires porting the function to Py2HWSW's core interface first.

\texttt{CreateVersatClass} has the following signature:

\begin{lstlisting}[caption={iob\_versat.py}]
CreateVersatClass(versat_spec, versat_top, versat_extra, build_dir,
                  axi_data_w, debug_path=None, profile=None, extraFlags=None)
\end{lstlisting}

It takes the specification file, the top-level unit name, an optional extra units folder, the SoC build directory, and the AXI data width. It then:

\begin{itemize}
  \itemsep-0.5em
  \item runs the \texttt{versat} executable found on \texttt{PATH}, not \texttt{./versat}, so the repository root, where \texttt{make -j versat} places the binary, must be on \texttt{PATH}. It always passes \texttt{-{}-debug}, \texttt{-d} to enable the DMA, \texttt{-b} set to the AXI data width, \texttt{-t} set to the top-level unit, \texttt{-o <build\_dir>/hardware/src}, and \texttt{-O <build\_dir>/software}, and adds \texttt{-u}, \texttt{-g}, and \texttt{-{}-profile} when \texttt{versat\_extra}, \texttt{debug\_path}, and \texttt{profile} are given; \texttt{extraFlags}, when given, is appended verbatim as one final raw argument to the \texttt{versat} invocation;
  \item saves Versat's output to \texttt{<build\_dir>/software/versatSetup.txt} and, whenever that file exists, reuses it instead of running Versat again, so the file must be deleted for a changed specification to take effect;
  \item reads the \texttt{ADDR\_W} line of that output and declares a \texttt{MAX\_CONFIG} register near the top of that address range, so that the SoC reserves the accelerator's whole address space;
  \item returns an \texttt{iob\_versat} class with clock, clock-enable, and reset signals and an IOb native CPU interface. Because the DMA enabled by \texttt{-d} always counts as a databus user, Versat always reports \texttt{HAS\_AXI}, so the class also always declares the \texttt{USE\_EXTMEM} configuration and an AXI master port that connects the accelerator directly to external memory.
\end{itemize}

The \href{https://github.com/IObundle/iob-soc-versat}{iob-soc-versat} repository is the reference integration. Its SoC class, which inherits from IOb-SoC's \texttt{iob\_soc}, integrates Versat in two class methods. \texttt{\_create\_submodules\_list} creates the accelerator class and lists it among the SoC's submodules:

\begin{lstlisting}[caption={iob\_soc\_versat.py}]
cls.versat_type = CreateVersatClass(
    pc_emul,
    VERSAT_SPEC,
    GetTestName(),
    VERSAT_EXTRA_UNITS,
    GetBuildDir("iob_soc_versat"),
    os.path.realpath(os.path.join(cls.setup_dir, "../debug/")),
)

super()._create_submodules_list(
    [
        # ... bus interfaces ...
        iob_vexriscv,
        cls.versat_type,
        iob_reset_sync,
    ]
    + extra_submodules
)
\end{lstlisting}

\texttt{\_create\_instances} then instantiates it as a peripheral named \texttt{VERSAT0}:

\begin{lstlisting}[caption={iob\_soc\_versat.py}]
if cls.versat_type in cls.submodule_list:
    cls.versat = cls.versat_type("VERSAT0", "Versat accelerator")
    cls.peripherals.append(cls.versat)
\end{lstlisting}

iob-soc-versat uses a different revision of iob-versat, whose \texttt{CreateVersatClass} takes an extra leading \texttt{pc\_emul} argument. With the signature in this repository, that argument is dropped and the AXI data width follows the build directory. Once the SoC is set up, its firmware includes the generated \texttt{versat\_accel.h} and calls \texttt{versat\_init} with the base address the SoC assigns to the \texttt{VERSAT0} peripheral. The repository's Makefile targets each enter a \texttt{nix-shell} environment defined by its own \texttt{default.nix}. This repository's \texttt{versat.nix} packages Versat itself.
